线性模型对整个输入空间只承诺一个函数形状，系数一旦定下，每个特征的作用在哪里都一样。树和邻域方法放弃了这份承诺：它们认为不同区域可以适用不同规则。放弃换来了灵活，也换来了三个新问题——空间该怎么划分，哪些点算彼此相近，以及一小块局部里的证据够不够支撑一个结论。

这三个问题的答案都是几何的。决策树用轴对齐的切分把空间剁成矩形盒子，于是交互和阈值效应不用显式写出来就被吃下；相应地，整套假设绑死在坐标轴上，换一组基就换一个模型。集成不去修单棵树，而是同时养很多棵：Bagging 与随机森林靠重采样和特征子采样制造分歧，用平均削掉方差；Boosting 让每棵新树去拟合当前损失的负梯度，一小步一小步把偏差磨掉。KNN 最直接，它连模型都懒得建，把``相近''交给一个距离函数，然后暴露出两个躲不掉的追问：尺子是谁定的，以及在这个维度下邻域还算不算局部。

\begin{center}
\begin{tikzpicture}[>=stealth]
% ---- 树：轴对齐盒子 ----
\draw[thick] (0,0) rectangle (2.2,2.2);
\draw (1.1,0) -- (1.1,2.2);
\draw (0,1.4) -- (1.1,1.4);
\draw (1.1,0.8) -- (2.2,0.8);
\node at (1.1,-0.45) {\footnotesize 树：矩形盒子};
% ---- 集成：多套分区叠加 ----
\begin{scope}[shift={(3.4,0)}]
\draw[thick] (0,0) rectangle (2.2,2.2);
\draw[gray] (0.7,0) -- (0.7,2.2);
\draw[gray] (0.7,1.6) -- (2.2,1.6);
\draw[gray] (1.5,0) -- (1.5,1.6);
\draw[gray] (0,0.6) -- (2.2,0.6);
\draw[gray] (1.9,0.6) -- (1.9,2.2);
\draw[gray] (0,1.9) -- (1.9,1.9);
\node at (1.1,-0.45) {\footnotesize 集成：多套分区平均};
\end{scope}
% ---- KNN：查询点周围的球 ----
\begin{scope}[shift={(6.8,0)}]
\draw[thick] (0,0) rectangle (2.2,2.2);
\foreach \p in {(0.4,0.5),(0.9,1.5),(1.3,1.1),(1.6,1.7),(0.6,1.9),(1.9,0.7),(1.2,0.4),(1.7,1.3),(0.3,1.2),(2.0,1.9)}
  \draw \p circle (2.2pt);
\draw[dashed] (1.3,1.3) circle (0.62);
\draw[very thick] (1.15,1.15) -- (1.45,1.45);
\draw[very thick] (1.45,1.15) -- (1.15,1.45);
\node at (1.1,-0.45) {\footnotesize KNN：查询点周围的球};
\end{scope}
\end{tikzpicture}
\end{center}

\emph{三种方法各自用一种形状回答``附近是什么''，后面的所有优点和失效都从这个形状长出来。}

图里的三种形状值得先记住，因为本章的失效模式全部可以追回到它。矩形盒子怕旋转，斜着的边界要用一串台阶去拼；一堆盒子取平均能把抖动压小，却挪不动它们共同的偏心；而球形邻域在维度升高时会膨胀到装下半个空间，此时最近与最远的差距缩进噪声里，``最近''不再携带信息。

本章的数学工具因此不是同一套可微目标函数。不纯度下降与加权方差减少解释切分怎么选，方差---相关性分解解释平均能省下多少，函数空间上的梯度下降解释 Boosting 每一步在干什么，而距离的集中现象解释 KNN 在高维为什么失灵。把它们硬拗成``都是梯度下降''会丢掉这些方法各自的诊断线索。

\subsection{怎么读这一章}

快读时给三篇各留一个判断：树的每一刀都是当下最优的贪心选择，且这个选择依赖于你用的坐标系；集成里 Bagging 动方差、Boosting 动偏差，因此失效的样子也不同；KNN 把度量和局部样本量直接写成了模型假设。第一遍不必记住 Gini 和熵的差别，也不必背下 boosting 的各种变体，但要能说清楚树深、树间相关和距离集中这三件事各自会怎样把方法拖垮。

若你带着具体问题进来，可以按症状定位：模型在验证集上抖得厉害，看第一篇的深度控制和第二篇的方差削减；加了很多树却不再变好，看第二篇关于相关性地板与偏差的讨论；特征重要性的排名和领域知识对不上，第二篇有一节专讲它偏在哪；特征维数几百而样本只有几千，第三篇的两节高维分析是必读的。

\subsection{本章篇目}

以下篇目按概念依赖排列；若从具体问题进入，也可以先打开最接近当前困惑的一篇，再沿篇内链接回补。

\begin{itemize}
\tightlist
\item \hyperref[sec:13-Machine-Learning/03-Trees-Ensembles-and-Neighbors/01]{``决策树把预测写成递归分区''}；
\item \hyperref[sec:13-Machine-Learning/03-Trees-Ensembles-and-Neighbors/02]{``集成用平均与逐步加法改变误差''}；
\item \hyperref[sec:13-Machine-Learning/03-Trees-Ensembles-and-Neighbors/03]{``KNN 把度量直接变成归纳偏置''}。
\end{itemize}

三篇读完，你手上会多出一组不依赖可微性的建模手段，以及一份关于它们何时不可信的清单。再往后走，当划分不再由标签驱动、连``分对了没有''都失去外部裁判时，问题会变成另一种形态——那是聚类要面对的局面。
